\documentclass[12pt, a4paper]{article}
\usepackage[margin=1in]{geometry}
\usepackage{hyperref}
\usepackage{graphicx}
\usepackage{amsmath}
\usepackage{amssymb}
\usepackage{cite}
\usepackage{fancyhdr}
\usepackage{xcolor}
\usepackage{listings}
\usepackage[T1]{fontenc}
\usepackage{setspace}

% Fix fancyhdr headheight warning
\setlength{\headheight}{14.5pt}

% Define colors
\definecolor{darkblue}{RGB}{0, 51, 102}
\definecolor{lightgray}{RGB}{240, 240, 240}

% Header and footer
\pagestyle{fancy}
\fancyhf{}
\rhead{\textit{Research Proposal}}
\lhead{\textit{Autism Teaching AI}}
\cfoot{\thepage}

% Title styling
\title{
    \textbf{\color{darkblue}Adaptive Multimodal AI Teaching System} \\
    \textbf{\color{darkblue}for Autism Intervention}
    \\[0.5in]
    \large{Research Proposal}
}
\author{}
\date{\today}

\begin{document}

\maketitle

\begin{center}
    \textit{PhD Research Proposal} \\
    Department of Computer Science and Engineering \\
    [Your Institution]
\end{center}

\newpage
\tableofcontents
\newpage

% ============================================================================
\section{Executive Summary}
% ============================================================================

This proposal outlines a novel approach to developing an \textbf{AI-powered adaptive teaching system} designed to assist in autism spectrum disorder (ASD) intervention. The system combines Vision-Language Models (VLAs), multimodal behavioral analysis, and dynamic dialogue generation to create a personalized, scalable teaching assistant that learns from therapist interactions and adapts to individual student needs in real-time.

\textbf{Key Innovation:} Unlike existing therapy robots (which are pre-scripted or presentation-only), our system performs \textit{real-time behavioral observation} of students and \textit{dynamically adapts teaching strategies} based on detected engagement, confusion, and frustration levels.

\textbf{Implementation Flexibility:} The core software is platform-agnostic. It can be deployed to any robot platform available in the lab (Pepper, NAO, TurtleBot, custom platforms, or even simulation for initial phases).

\textbf{Research Timeline:} 
\begin{itemize}
    \item \textbf{Phase 1 (8 weeks):} Software architecture + behavior analysis pipeline on public datasets
    \item \textbf{Phase 2 (Months 3-6):} Integration with real therapy data + adaptation
    \item \textbf{Phase 3 (Months 7-12):} Deployment and feasibility evaluation
\end{itemize}

\textbf{Expected Outcomes:}
\begin{enumerate}
    \item Behavioral prediction model achieving $\geq 80\%$ accuracy on autism-specific behaviors
    \item Fine-tuned Vision-Language Model for adaptive educational interaction
    \item End-to-end teachable AI system architecture
    \item Feasibility insights from evaluation with real students
    \item 2-3 peer-reviewed publications (AI + special education venues)
\end{enumerate}

\newpage

% ============================================================================
\section{Background and Motivation}
% ============================================================================

\subsection{Current State of Autism Intervention}

Autism spectrum disorder (ASD) affects over 1 in 36 children globally. Effective intervention requires intensive, personalized, multi-modal teaching. However:

\begin{itemize}
    \item \textbf{Limited therapist capacity:} Insufficient trained specialists to meet demand
    \item \textbf{High variability:} Teaching methods are not standardized; each therapist uses different approaches
    \item \textbf{Rigid educational technology:} Existing therapy robots are pre-scripted and cannot adapt to student behavior in real-time
    \item \textbf{Lack of measurement:} Limited automated real-time feedback on student engagement and learning progress
\end{itemize}

\subsection{Why Vision-Language Models Now?}

Recent advances in VLAs (UniVLA 2025, OpenVLA 2024, LLaRA 2025) enable:
\begin{itemize}
    \item End-to-end learning from video demonstrations (perfect for learning from therapy recordings)
    \item Parameter-efficient adaptation via LoRA (fine-tune on limited data)
    \item Simultaneous perception (visual observation of student) + reasoning (teaching strategy)
    \item Real-time multi-modal understanding (video + audio combined)
\end{itemize}

\textbf{This creates a unique opportunity:} Combine VLAs with behavioral analysis to create AI that truly understands \textit{how and why} a therapist teaches effectively, then replicate that approach.

% ============================================================================
\section{Research Objectives and Contributions}
% ============================================================================

\subsection{Primary Research Questions}

\begin{enumerate}
    \item Can we accurately predict student engagement/confusion from multi-modal therapy session recordings?
    \item Can real-time behavioral observation enable AI systems to adapt teaching strategies dynamically?
    \item How well can AI learn the \textit{reasoning} behind effective teaching by analyzing therapist-student interactions?
    \item What is the practical feasibility of deploying such a system in a therapeutic or educational setting?
\end{enumerate}

\subsection{Novel Contributions}

\begin{enumerate}
    \item \textbf{First integration:} Real-time behavioral perception + adaptive teaching strategy in autism education (not previously published)
    \item \textbf{Novel fine-tuning methodology:} Parameter-efficient VLA adaptation for therapeutic/educational interaction (distinct from manipulation/locomotion tasks)
    \item \textbf{Teaching pattern extraction:} Framework for analyzing therapy recordings to extract implicit teaching patterns via LLM + behavioral analysis
    \item \textbf{Platform-agnostic architecture:} Software designed to work with any robot platform, not tied to specific hardware
    \item \textbf{Feasibility insights:} Real-world validation approach for AI tutoring systems in special education
\end{enumerate}

% ============================================================================
\section{Technical Approach}
% ============================================================================

\subsection{System Architecture}

The system has four core AI components that connect observation to adaptive response:

\begin{equation}
\text{Video/Audio} \rightarrow \text{Behavior Analysis} \rightarrow \text{Strategy Selection} \rightarrow \text{Dialogue Generation}
\end{equation}

Robot platform deployment is optional based on lab resources. The software pipeline is the core thesis contribution.

\subsubsection{Component 1: Multimodal Behavioral Analysis}

\textbf{Input:} Video stream (student + therapist) + audio (speech + prosody)

\textbf{Processing:}
\begin{itemize}
    \item \textbf{Visual features:} Face landmarks + pose estimation (MediaPipe) + expression classification (fine-tuned LLaVA)
    \item \textbf{Audio features:} Speech transcription (Whisper) + prosody analysis (speech rate, pitch, energy via librosa)
    \item \textbf{Fusion:} Multi-modal model predicting student state $s_t$
\end{itemize}

\textbf{Output:} 
\begin{equation}
s_t = \arg\max_s P(s | V_t, A_t; \theta) \quad \text{where } s \in \{\text{engaged, confused, frustrated, disinterested}\}
\end{equation}

\textbf{Accuracy Target:} $\geq 80\%$ vs. therapist ground-truth annotations

\subsubsection{Component 2: Teaching Strategy Selector}

\textbf{Input:} Student state $(s_t)$, current topic/lesson, interaction history

\textbf{Processing:}
\begin{itemize}
    \item Extract teaching patterns from therapy transcripts using LLM analysis
    \item Build decision policy: \textit{``When student shows state $s$, use strategy $a$''}
    \item Strategies include: clarify concept, simplify language, use visual aids, change modality, take break, ask guiding question
\end{itemize}

\textbf{Output:} Selected teaching strategy + parameters

\subsubsection{Component 3: Dialogue Generation}

\textbf{Input:} Selected strategy, student state, current topic, interaction context

\textbf{Processing:}
\begin{itemize}
    \item LLM-based generation (GPT-3.5/4 or open-source Mixtral-8x7B)
    \item Prompt: \textit{``I am an autism specialist. This student is [state]. Teach [topic] using [strategy]. Be clear and patient.''}
    \item Optional TTS (Eleven Labs, glow-tts, or lab's existing system)
\end{itemize}

\textbf{Output:} Natural language teaching response (+ audio if TTS available)

\subsubsection{Component 4: Vision-Language Model (Optional Robot Grounding)}

\textbf{If deploying to a robot platform:}
\begin{itemize}
    \item Base pre-trained VLA (OpenVLA-7B or UniVLA)
    \item Fine-tune with LoRA on therapy interaction datasets
    \item Learn robot gestures, timing, and interaction patterns from therapist behavior
    \item Real-time perception-action loop
\end{itemize}

\textbf{Important:} This component only activates if/when robot deployment occurs. Phase 1-2 focus on the perception + strategy + dialogue pipeline.

\subsection{Technology Stack}

\begin{table}[h]
\centering
\begin{tabular}{|l|l|}
\hline
\textbf{Component} & \textbf{Technology/Framework} \\
\hline
Behavioral Analysis & LLaVA, Qwen-VL (Hugging Face) \\
\hline
Pose + Face Detection & MediaPipe \\
\hline
Audio Analysis & Whisper (speech-to-text) + librosa (prosody) \\
\hline
Strategy Selection & Decision tree / learned policy (numpy/sklearn) \\
\hline
Dialogue Generation & GPT-3.5/4 API or Mixtral-8x7B \\
\hline
Optional: VLA Fine-tuning & OpenVLA-7B + LoRA (PyTorch) \\
\hline
Optional: TTS & Eleven Labs API or open-source TTS \\
\hline
Optional: Robot Deployment & ROS2, Pepper SDK, NAO SDK, or platform of choice \\
\hline
\end{tabular}
\end{table}

All technologies are either open-source or have free tier APIs.

% ============================================================================
\section{Datasets and Data Resources}
% ============================================================================

\subsection{Phase 1 (Demo): Public Datasets}

\begin{itemize}
    \item \textbf{MMASD:} MultiModal ASD Dataset (6+ hours, public therapy sessions)
    \item \textbf{Expression recognition:} FER2013, AffectNet (general expressions)
    \item \textbf{Educational videos:} Khan Academy, OpenStax, TED-Ed (for strategy synthesis)
\end{itemize}

No permissions needed; freely available.

\subsection{Phase 2+ (Full Implementation): Real Therapy Data}

\begin{itemize}
    \item Partner with autism intervention centers for therapy video/audio
    \item Target: 50-100 hours (can start smaller: 10-20 hours)
    \item De-identify + HIPAA compliance
    \item Therapist annotations for ground-truth behavioral labels
\end{itemize}

\subsection{Data Privacy}

\begin{itemize}
    \item All public datasets in Phase 1; no patient data required
    \item Phase 2 therapy data: de-identified, HIPAA-compliant, IRB-approved
    \item Models released anonymously; real patient data never shared
\end{itemize}

% ============================================================================
\section{Project Timeline and Milestones}
% ============================================================================

\subsection{Phase 1: Software Architecture Demo (Weeks 1-8)}

Build a working end-to-end AI teaching system with public datasets.

\begin{center}
\begin{tabular}{|p{2cm}|p{6.5cm}|p{2.5cm}|}
\hline
\textbf{Week(s)} & \textbf{Milestone} & \textbf{Deliverable} \\
\hline
1-3 & Behavioral analysis pipeline on MMASD & Model checkpoint + notebook \\
\hline
3-5 & Strategy extraction from transcripts & Decision policy document \\
\hline
5-8 & Dialogue generation + integration & End-to-end demo video \\
\hline
\textbf{Week 8} & \textbf{Demo Ready} & Reproducible code + short paper \\
\hline
\end{tabular}
\end{center}

\subsection{Phase 2: Real Data Integration (Months 3-6)}

\begin{enumerate}
    \item Request + receive therapy data from partners (HIPAA-compliant)
    \item Fine-tune behavioral model on real student interactions
    \item Adapt dialogue generation for real therapy context
    \item First publication: ``Real-Time Behavioral Adaptation in AI Teaching Systems for Autism''
\end{enumerate}

\textbf{Robot integration (if lab has platform available):} Optional during this phase. Can add VLA component and connect to existing robot.

\subsection{Phase 3: Evaluation and Feasibility (Months 7-12)}

\begin{enumerate}
    \item IRB approval for small feasibility study
    \item Deploy system (with or without robot, depending on lab resources)
    \item Evaluate with 5-10 students with autism
    \item Measure: engagement, learning outcomes, user acceptance
    \item Second publication + thesis
\end{enumerate}

% ============================================================================
\section{Evaluation Metrics}
% ============================================================================

\subsection{Phase 1: Software Quality}

\begin{itemize}
    \item Behavioral prediction accuracy: $\geq 80\%$ vs. therapist labels
    \item Strategy selection coherence: Expert (therapist) ratings 4+ / 5
    \item Dialogue naturalness: Human raters 4+ / 5 on Likert scale
    \item Code reproducibility: All code + models on GitHub/Hugging Face
\end{itemize}

\subsection{Phase 2: Real Data Performance}

\begin{itemize}
    \item Behavioral model validation: Inter-rater agreement (Cohen's $\kappa \geq 0.70$)
    \item Adaptation effectiveness: Student engagement before vs. after strategy selection
\end{itemize}

\subsection{Phase 3: Real-World Feasibility}

\begin{itemize}
    \item Student engagement: Measured via gaze, gesture, speech
    \item Learning outcomes: Pre/post assessment scores
    \item Therapist usability: System usability scale (SUS) score
    \item Feasibility: Can system be integrated into real therapy workflows?
\end{itemize}

% ============================================================================
\section{Related Work and Positioning}
% ============================================================================

\subsection{What Exists}

\begin{itemize}
    \item Therapy robots (Pepper, NAO): Pre-scripted, cannot adapt to student behavior
    \item Autism behavior recognition (Deng et al., 2024): Can detect behaviors but no teaching integration
    \item Vision-Language Models (OpenVLA, UniVLA): SOTA for robotics but not adapted to education
    \item LLM-based tutoring: Text-only, no visual perception of student
\end{itemize}

\subsection{What's Missing (Our Contribution)}

\begin{enumerate}
    \item \textbf{Real-time behavioral observation} informed by video analysis
    \item \textbf{Adaptive teaching strategies} based on observed student state
    \item \textbf{Learning from therapist behavior} via multimodal analysis
    \item \textbf{End-to-end system} connecting perception $\rightarrow$ reasoning $\rightarrow$ response
\end{enumerate}

This work is genuinely novel: No published paper combines all four elements.

% ============================================================================
\section{Estimated Resources}
% ============================================================================

\subsection{Computational}

\begin{itemize}
    \item GPU: 1x NVIDIA A100 (24GB) or similar for model training
    \item Storage: 500GB-1TB for datasets + checkpoints
    \item Cloud compute (optional): AWS/GCP for large-scale fine-tuning
\end{itemize}

\subsection{Data}

\begin{itemize}
    \item Public datasets (MMASD, FER2013): Free
    \item Real therapy data: Depends on partner availability (often shared for research)
\end{itemize}

\subsection{Software}

\begin{itemize}
    \item Open-source: PyTorch, Hugging Face, MediaPipe, librosa
    \item APIs (optional): GPT-3.5/4 (pay-as-you-go), Eleven Labs TTS (free tier exists)
    \item Robot platform: Use whatever lab already has (no new hardware purchase required)
\end{itemize}

\subsection{Estimated Budget (Optional)}

\begin{table}[h]
\centering
\begin{tabular}{|l|r|}
\hline
Cloud GPU compute (12 months) & \$2000-4000 \\
\hline
API costs (Whisper, GPT-4, TTS) & \$500-1000 \\
\hline
Robot platform & \$0 (use lab's existing resources) \\
\hline
\textbf{Total} & \textbf{\$2500-5000} \\
\hline
\end{tabular}
\end{table}

\textbf{Note:} Hardware costs are minimal because we use robots the lab already has.

% ============================================================================
\section{Risks and Mitigation}
% ============================================================================

\begin{table}[h]
\centering
\begin{tabular}{|p{2.5cm}|p{3.5cm}|p{1.5cm}|}
\hline
\textbf{Risk} & \textbf{Mitigation} & \textbf{Likelihood} \\
\hline
Therapy data unavailable & Start with MMASD public data; frameworks work with any dataset & Low \\
\hline
VLA adaptation underperforms & Use LLM dialogue as backup; combine approaches & Low \\
\hline
IRB approval delays & Phase 1-2 complete before human studies needed & Low \\
\hline
Robot integration issues & Software works without robot; add later if needed & Medium \\
\hline
\end{tabular}
\end{table}

% ============================================================================
\section{Conclusion}
% ============================================================================

This proposal outlines a feasible, high-impact PhD research project combining recent advances in Vision-Language Models, multimodal analysis, and adaptive dialogue to address a genuine gap in autism intervention technology.

\textbf{Key strengths:}
\begin{itemize}
    \item \textbf{Novel:} First system to combine behavioral observation + adaptive teaching in autism context
    \item \textbf{Feasible:} Uses proven technologies (OpenVLA, LLaVA, Whisper); no fundamental research blockers
    \item \textbf{Scalable:} Software-first design works with any robot platform (or none)
    \item \textbf{Impactful:} Addresses real need in special education and autism intervention
    \item \textbf{Publishable:} Clear path to 2-3 high-quality publications
\end{itemize}

Phase 1 (proof-of-concept) can begin immediately with public datasets.

\newpage

% ============================================================================
\section*{References}
% ============================================================================

\begin{thebibliography}{99}

\bibitem{Mishra2024} Mishra, A., et al. (2024). Human-mediated Large Language Models for Robotic Intervention in Children with Autism Spectrum Disorders. \textit{arXiv:2402.00260}.

\bibitem{Kim2024} Kim, A., et al. (2024). OpenVLA: Towards a Foundation Model for Vision and Action. In \textit{ICML 2024}.

\bibitem{Bu2025} Bu, Z., et al. (2025). UniVLA: A Unified Vision-Language Action Model. \textit{RSS 2025}.

\bibitem{Deng2024} Deng, J., et al. (2024). Audio-Visual Autism Behavior Recognition. \textit{arXiv:2406.02554}.

\end{thebibliography}

\end{document}